\documentclass[10pt]{article}
\usepackage[margin=0.7in]{geometry}
\usepackage[T1]{fontenc}
\usepackage{lmodern}
\usepackage{enumitem}
\setlist[enumerate]{leftmargin=*, itemsep=0.65em}
\setlist[enumerate,2]{label=(\Alph*), leftmargin=2em, itemsep=0.05em, topsep=0.25em}
\pagestyle{plain}

\begin{document}

\begin{center}
  {\large\bfseries Video 1 Quiz: Propositional Logic}\\[0.25em]
  \textit{Choose the best answer for each question.}
\end{center}

\noindent Name:\ \rule{3.2in}{0.4pt}\hfill Date:\ \rule{1.25in}{0.4pt}

\vspace{0.7em}

\begin{enumerate}
  \item Which example sentence is used in the video's introduction to statements?
  \begin{enumerate}
    \item The square root of $3$ is irrational.
    \item The real number $\sqrt{2}$ is rational.
    \item Every prime greater than $2$ is odd.
    \item The angles of every triangle sum to $180$ degrees.
  \end{enumerate}

  \item In the video, what is the name for the proposition ``Duke University is in Durham, North Carolina, or all real numbers are rational''?
  \begin{enumerate}
    \item A conjunction.
    \item A disjunction.
    \item A conditional.
    \item A biconditional.
  \end{enumerate}

  \item In the truth table for the conditional ``if $P$, then $Q$,'' which case makes the conditional false?
  \begin{enumerate}
    \item $P$ is true and $Q$ is false.
    \item $P$ is false and $Q$ is true.
    \item $P$ is false and $Q$ is false.
    \item $P$ is true and $Q$ is true.
  \end{enumerate}

  \item What does the double-arrow implication $R \Rightarrow S$ assert in the video?
  \begin{enumerate}
    \item $R$ and $S$ are individual propositions with opposite truth values.
    \item The ordinary conditional ``if $R$, then $S$'' is a tautology.
    \item $R$ is true only when $S$ is false.
    \item $R$ and $S$ are the same letter.
  \end{enumerate}
\end{enumerate}

\vspace{0.7em}
\noindent\textit{This quiz assesses the main ideas from the propositional-logic video.}

\end{document}
